\begin{lemma}
Let $S$ be a $3$-uniform $3$-simple local extremal setup at scale $D$
for the indexed hypergraph $\mathcal F$ and edge index $f$.
Let $N$ be the set of edge indices other than $f$ meeting $f$ and put
$q(g)=|\{h\in N:\mathcal F(g)\cap\mathcal F(h)=\varnothing\}|$.
Then $T_2\le\frac1{12}\sum_{g\in V_2}q(g)^2$.
\end{lemma}
\begin{proof}
Use the data of \noderef{ThreeUniformThreeSimpleLocalSetup}.
On the finite vertex type consisting of indices in $V_2$, form the graph
whose adjacency is disjointness of the two indexed hyperedges.
This is symmetric and irreflexive: every hyperedge has size three and
so is nonempty. Its triangles correspond bijectively to the three-element
subsets of $V_2$ consisting of pairwise disjoint hyperedges, by forgetting
the membership proofs of the subtype. Conversely each such subset lifts
uniquely to the subtype. The defining formula for $T_2$ counts exactly
these subsets, since $V_2\subseteq N$ and nonadjacency in the line graph
for distinct indices means disjointness.

Each $g\in V_2$ meets $f$ in exactly one vertex, by the edge-shape clause
of the setup. Label $g$ by that vertex, identifying the three vertices
of $f$ with $\{0,1,2\}$. Adjacent vertices cannot have the same label,
since that would give a common vertex of their hyperedges. Thus the graph
is tripartite. By \noderef{TriangleCountTripartiteCherries}, its triangle
count is at most one twelfth of the sum of its squared degrees.
For each $g$, forgetting subtype membership injects its neighbor set
into $\{h\in N:\mathcal F(g)\cap\mathcal F(h)=\varnothing\}$.
Its degree is therefore at most $q(g)$. Both counts are nonnegative,
so squaring and summing preserves this inequality and proves the claim.
\end{proof}
